% ============================================================================
%  1. INTRODUCTION
%
%  A&A exige que el manuscrito empiece por "1. Introduction". La numeración la
%  lleva LaTeX: nunca escribir el número a mano.
%
%  Jerarquía disponible:
%      \section{Título}        -> 1.
%      \subsection{Título}     -> 1.1
%      \subsubsection{Título}  -> 1.1.1
%      \paragraph{Título}
%  Toda sección necesita un título corto y descriptivo.
%
%  Cross-referencing: \label en TODA sección, figura, tabla y ecuación
%  referenciada; \ref para citarlas. A&A lo convierte en hipervínculos HTML.
%  Atajos de aapaper.sty: \secref{} \figref{} \tabref{} \eqnref{} \appref{}
%
%  Objetos astronómicos: envolverlos SIEMPRE en \object{}, también dentro de
%  tablas pequeñas. Genera el archivo .obj que enlaza con SIMBAD/ALADIN.
%      \object{Sgr A*}   \object{S2}
%
%  Abreviaturas: expandir la primera vez que aparecen en el texto principal.
% ============================================================================

\section{Introduction}
\label{sec:introduction}

Texto de la introducción \citep{baker66}. Esa cita es un placeholder: un
\texttt{thebibliography} sin ninguna entrada es un error de LaTeX, así que el
esqueleto trae una para compilar. Bórrala junto con la entrada de
\texttt{refs.bib} en cuanto tengas tus propias referencias.

% ----------------------------------------------------------------------------
%  Ejemplos de uso — descomentar y adaptar
% ----------------------------------------------------------------------------
%
%  Citas (natbib con la puntuación A&A ya configurada en main.tex):
%      \citet{gillessen17}          -> Gillessen et al. (2017)
%      \citep{gillessen17}          -> (Gillessen et al. 2017)
%      \citep[see][chap.~2]{do19}   -> (see Do et al. 2019, chap. 2)
%      \citep{a18,b19}              -> (A et al. 2018; B et al. 2019)
%
%  Referencias cruzadas:
%      as shown in \secref{sec:methods}
%      \figref{fig:orbit} displays the projected orbit
%      \tabref{tab:elements} lists the orbital elements
%      from \eqnref{eq:geodesic}
%
%  Objetos:
%      The star \object{S2} orbits \object{Sgr A*} with a period of 16 yr.
%
%  Tipografía A&A (ver el checklist completo en la skill aa-paper):
%      20\,000 km              miles con espacio fino
%      1950--1985              rango con en-dash
%      $-30$~K                 menos matemático + espacio fijo
%      $v = 7650$~km~s$^{-1}$  unidad con índice negativo, nunca km/s
%      $T_\mathrm{eff}$        subíndice-etiqueta en redonda
%      \ion{H}{ii}             ion, romano en minúsculas
% ----------------------------------------------------------------------------
